\chapter{Electronics and Wiring}

\section{Overview}

The electronics subsystem includes the STM32 Nucleo-L010RB board, stepper motor drivers, stepper motors, power supply, and connections to the SDRs and host computer.

\section{Wiring Guidelines}

Use the wiring diagram shown in Fig. \ref{fig:wiring} to connect the components of the system.
The pins on the STM32 are shown with two labels denoting the pin numbers of the STM32, and the pins
of the corresponding Arduino header that match and are shown on the board.

\begin{figure}[h!]
    \footnotesize
    \begin{centering}
    \includesvg[width=\textwidth]{figs/wiring}
    \end{centering}
    \caption{Wiring Diagram for STM32 and Motor Drivers.}
    \label{fig:wiring}
\end{figure}

A +12V power supply is required to power the motor drivers and is omitted here for brevity.

\section{Stepper Motor Drivers}

Each stepper driver (TB6600) connects to:
\begin{itemize}
  \item A Stepper Motor
  \item Power: +12~V and GND.
  \item Control signals from STM32: PUL, DIR, and ENABLE.
\end{itemize}

Unless another pulse/revolution ratio is desired and sent to the STM32 via the \hyperref[sec:serial_protocol]{serial protocol},
set the TB6600 motor drivers to 800 pulse/revolution.


\section{STM32 Nucleo Connections}

\begin{itemize}
  \item USB connection to the host computer for programming and serial communication.
  \item GPIO pins connected to each driver's PUL(-), DIR(-), and EN(-) pins.
  \item 5V pin connected to each driver's PUL(+), DIR(+), and EN(+) pins.
\end{itemize}

\section{SDR Connections}

\begin{itemize}
  \item The transmit SDR \& receiver SDR connects to the host computer via USB.
  \item The transmit SDR connects to WIFI band parabolic antenaa.
  \item The receive SDR connects to the AUT.
\end{itemize}

\begin{figure}[h!]
    \footnotesize
    \begin{centering}
    \includesvg[width=\textwidth]{figs/wiring_computer}
    \end{centering}
    \caption{Wiring Diagram for STM32 and Motor Drivers.}
    \label{fig:wiring}
\end{figure}
